% ===== =======================================================
% Preambulo
% ============================================================

% Tipo de documentos
\documentclass[11pt,a4paper]{article}

% Rutas relativas: no dependen del usuario ni de la ubicación del repositorio.
% Detectar el directorio desde el que TeXstudio o el compilador ejecutan Master.tex.
\usepackage[strings]{underscore}
\newcommand{\joinpath}[2]{#1/#2}
\IfFileExists{../B_Packs/Packs.tex}{%
	% Compilación desde D_LaTeX/A_Master (TeXstudio).
	\newcommand{\basepath}{../..}
	\newcommand{\latexpath}{..}
}{%
	\IfFileExists{D_LaTeX/B_Packs/Packs.tex}{%
		% Compilación desde TP_01.
		\newcommand{\basepath}{.}
		\newcommand{\latexpath}{D_LaTeX}
	}{%
		\IfFileExists{TP_01/D_LaTeX/B_Packs/Packs.tex}{%
			% Compilación desde la raíz del repositorio.
			\newcommand{\basepath}{TP_01}
			\newcommand{\latexpath}{TP_01/D_LaTeX}
		}{%
			\PackageError{Master}{No se encontro TP_01/D_LaTeX}{Abrir Master.tex desde el repositorio.}%
		}%
	}%
}
\newcommand{\outputspath}{\joinpath{\basepath}{B_Outputs}}

% Cargar paquetes
\input{\joinpath{\latexpath}{B_Packs/Packs.tex}}



% ============================================================
% Cuerpo
% ============================================================
\begin{document}
	%-------------------------------------------------------------
	% Portada
	%-------------------------------------------------------------	
	{\footnotesize Big Data --- Maestría en Economía, UdeSA\par}
	\vspace{0.5em}
	{\LARGE\bfseries Trabajo Práctico 01\par}
	\vspace{0.15em}
	{\large Antequera Álvaro, Fariña Valentina, Vicencio Joel\par}
	\vspace{0.7em}
	\hrule height 0.6pt
	\vspace{0.5em}
	%-------------------------------------------------------------
	% Secciones
	%-------------------------------------------------------------
	\section*{Declaración de uso de IA}
	En el presente trabajo se utilizaron herramientas de inteligencia artificial para la elaboración y revisión del código de Python, la organización de la estructura del documento en LaTeX y la incorporación de tablas, figuras y ecuaciones. También se utilizaron para la búsqueda y revisión de bibliografía académica que permitiera fundamentar la selección de algunas de las variables explicativas incluidas en la Parte I.
	
	\input{\joinpath{\latexpath}{C_Section/parte_01_v2.tex}}
	
	\input{\joinpath{\latexpath}{C_Section/parte_02.tex}}
	
	\input{\joinpath{\latexpath}{C_Section/parte_03.tex}}
	
	\input{\joinpath{\latexpath}{C_Section/parte_04.tex}}
	
	
	%-------------------------------------------------------------
	% Bibliografia
	%-------------------------------------------------------------
	
	\nocite{*}
	
	\bibliographystyle{apalike}
	\bibliography{\joinpath{\latexpath}{D_Bibliografia/biblio.bib}}
\end{document}